\section{April Showers 2013}

April showers have chased me indoors. Last night I made enough money playing \textbf{Texas Hold'em} to cover the purchase of the UR/KRUR books and Della Riveria's \emph{Magic World of the Heros}, so you guys are off the hook. But thanks for the offers. Don't knock playing cards which are clearly of Hermetic origin, when you consider how perfectly organized they are. Look at how many various games can be derived from them. Video games, on the other hand, fall into just a few types: shooters, role playing, and maybe another I can't think of. They are made to create fast-reacting zombies. With poker you learn to accept the delicate balance between skill and luck, while learning to discern men's wills. Plus skill at cards and similar amusements will prepare you for success at courtly life.

I have been planning another post on Darwinism, but I just found one by \textbf{Julius Evola} in the \emph{Introduction to Magic}. So, I'll start with that and perhaps add some comments. Unlike more recent attempts to challenge Darwinism on its own terms, such as Creation Science or Intelligent Design, Evola attacks it from a metaphysical viewpoint. If a biologists simply assumes only material processes are in play, then Darwinian evolution seems to account for the appearances. Evola does not dispute any biological facts; instead, he proposes that the idea of an involution also accounts for the same facts. I'll let the essay speak for itself for now. But to be prepared, read up on Aristotle and Aquinas, since Evola says they best express his views. Specifically, look up on Gornahoor about matter and form.

There are, of course, many other essays of interest, out of about 80 or so. Perhaps I'll take the time to create a table of contents and create a poll. Of course, someone else may decide to make that table of contents; they can be found on google books. I am also considering making further translations available only in pdf format; I am willing to consider any feedback on this.

I haven't forgotten Guido De Giorgio's section on the caste of workers. This group is given short shrift in writings on tradition, which are too often focused on spiritual and political issues. Arts, crafts, farming, and other commercial activities are crucial to any society.

I will also be providing some insights into the logic of human societies with historical examples. It should be interesting to see how we got to here from there.

I have some letters from the young Julius Evola to the great Italian philosopher \textbf{Giovanni Gentile} that I will make available. These will show the basis for Evola's own philosophy of magical idealism. I know later on Evola criticized Gentile, but a family dispute something different from a war against an enemy. That will be a prelude to a fuller exploration of Gentile's system of \textbf{Actualism}. Of course, Evola's own system of Magical Idealism is interesting in its own right, not least of all because he combines philosophical reasoning with traditional teachings. I'm afraid, however, that there are only 7 people who would be interested in it, and I can't be sure I am reaching all of them.

There is just one more chapter in part 3 of Sintesi that will be sent out to the mailing list of the Society of Independents \textbf{SI}. After that, I will upgrade the software to be a two-way discussion. There are many interesting subscribers of varied backgrounds who have developed skills in certain areas. I hope, initially, to model it on the \textbf{Gruppo di UR}. Participants can provide an essay on a topic of interest to them; think about what you may want to contribute. Of course, there is the matter of the Sintesi. Evola really skates over the various spiritual types without much detail. Perhaps some of these chapters can be filled out with more detail and examples. I would hope several readers are interested in that project. If that works out, I will make more sections available.


\flrightit{Posted on 2013-04-13 by Cologero}

\begin{center}* * *\end{center}

\begin{footnotesize}
\begin{sffamily}
\texttt{Jason-Adam on 2013-04-15 at 11:48 said:}

Is there a connection between being an aristocrat and being a gambler ? I too am a Texas Holdem player... ... ... ... ... ... .

Are any of Gentile's books in English ? I've never read anything by him but a lot about him and want to start.

\hfill

\texttt{Cologero on 2013-04-15 at 13:04 said:}

Amazon is giving away for free the kindle edition of Gentile's The Reform of Education. It was intended for elementary school teachers, although I doubt that many teachers could even comprehend it today.

\hfill

\texttt{Mercurius on 2013-04-15 at 18:10 said:}

I've uploaded ``Theory of Mind as Pure Act'' to Scribd:

\url{http://www.scribd.com/doc/136095606/Gentile-Theory-of-Mind-as-Pure-Act}


\hfill

\texttt{graham on 2013-04-16 at 22:44 said:}

I look forward to the chapter on workers. What proportion of those drawn to tradition today do you suppose are by nature workers? What does it mean if a worker is drawn to esoterism? What should be the worker's response to the mechanization of the crafts? What is his mode of realization? I hope he can shed some light on these and other questions.

\hfill

\texttt{Balder on 2022-04-13 at 09:13 said:}

Gambling is a form of astral intoxication and world glamour and can easily turn into a serious vice. The worst dregs in rehabilitation centres are gaming addicts. Card games in themselves are a good way to enjoy pasttime with friends, although I prefer chess.

I found Evola's Magical Idealism far too ``technical'' (for lack of a better word) to arouse my serious interests.

\end{sffamily}
\end{footnotesize}

